\chapter{局部紧空间上的积分.}

现代积分概念的发展，与函数观念的演进以及自 19 世纪初以来对实变量数值函数的深入研究紧密相连。众所周知，欧拉对函数概念的构想已经相当一般，因为在他看来，一条被每条平行于 Oy 轴的直线只交于一点的给定“任意”曲线就定义了一个函数 \(y = f(x)\)（参见第 196 页）；但是，与其同时代的大多数人一样，他拒绝承认这样的函数可以“解析地”表出。直到傅里叶的工作之前，这一观点几乎没有什么变化；而傅里叶所发现的把不连续函数表示为三角级数之和的可能性\(^{1}\)，对后世产生了决定性的影响。说实在的，傅里叶的证明完全缺乏严格性，其有效范围也未显现清楚；然而积分公式

\[
a _ {n} = \frac {1}{\pi} \int_ {- \pi} ^ {+ \pi} \phi (x) \cos n x d x, b _ {n} = \frac {1}{\pi} \int_ {- \pi} ^ {+ \pi} \phi (x) \sin n x d x (n \geq 1)\tag{22.1}
\]

它们给出函数 \(\phi\) 展开为傅里叶级数时的系数；只要假定 \(\phi\) 连续且逐段单调，这些公式便有一种直观上显而易见的意义。\(^{2}\)因此狄利克雷起初正是把自己局限于这些函数，在他建立傅里叶级数收敛性的那篇著名论文（{[}92{]}，第 I 卷，第 117–132 页）中即是如此；但在这项工作的末尾，他已经关心把他的结果推广到更广的函数类去了。众所周知，正是在这个场合，狄利克雷把傅里叶的观念精确化，定义了我们今天所理解的函数的一般概念；首先要澄清的一点自然是要知道，在哪些情形下仍能赋予公式 (22.1) 以意义。“当连续性{[}关于 \(\phi\) 的{]}的解有无穷多个\ldots{}”，狄利克雷说（同上出处，第 131–132 页），“那时就必须要求函数 \(\phi(x)\) 具有这样的性质：若以 \(a\) 和 \(b\) 表示 \(-\pi\) 与 \(+\pi\) 之间的两个任意量，则总能在 \(a\) 与 \(b\) 之间放上足够接近的其他量 \(r\) 与 \(s\)，使函数在从 \(r\) 到 \(s\) 的区间上保持连续。这一限制的必要性可以通过如下考虑来体会：级数{[}傅里叶级数{]}的各项都是定积分，而且要回到积分的基本概念。那时将会看到，函数的积分只在函数满足前述条件时才有意义。”

用现代的话说，狄利克雷似乎认为可积性等价于不连续点构成一个稀疏集这一事实；此外，在几行之后，他还指出了那个著名的例子：函数在 x 为有理数时等于 c，在 x 为无理数时取另一个不同的值 d，并断言这个函数“不能被代入”积分。他还宣布要在这一题目上继续工作，但这项工作从未发表，\(^{3}\)而且在此后 25 年间，似乎没有人试图沿这条道路前进，也许是因为在当时看来，考虑这类“病态”函数完全缺乏兴趣；无论如何，当黎曼于 1854 年（{[}259 a{]}，第 227–264 页）重新研究这个问题时（仍然与三角级数相关联\(^{4}\)），他感到有必要为他的工作辩护：“无论我们对力和物质的状态在无限小尺度下随时间与地点变化的方式如何无知，我们仍可确信，狄利克雷的研究所不适用的那些函数并不出现在自然现象中。然而”，他继续写道，“狄利克雷未处理的这些情形值得注意，理由有二。其一，正如狄利克雷本人在其工作末尾所指出的，这个题目与无穷小演算的原理有着极为密切的关系，可以用来给这些原理带来更多的清晰性和确定性。从这个观点看，对它的研究有直接的兴趣。其二，傅里叶级数的应用并不限于物理学中的研究；目前它们在纯粹数学的一个领域即数论中也得到了成功的应用，而在那里，重要的似乎恰恰是那些其三角级数展开尚未被狄利克雷研究过的函数”（{[}259 a{]}，第 237–238 页）。

黎曼的想法是从柯西使之恢复声誉的那个积分近似程序出发，确定一个函数 f 在有界区间 \([a, b]\) 上的“黎曼和”何时趋于一个极限（细分中各区间长度之最大者趋于 0）；他不费力地得到了这个问题如下形式的解：对每个 \(\alpha > 0\)，都存在 \([a, b]\) 的一个分成最大长度足够小的部分区间的分割，使得该分割中 f 的振幅 \(>\alpha\) 的那些区间的长度之和任意小。他还证明，这一条件不仅为连续且逐段单调的函数所满足，而且为可以有处处稠密的不连续点集的函数所满足。\(^{5}\)

黎曼的这篇论文直到他去世后于 1867 年才发表。但这一次，时代对这类研究已较为有利，“黎曼积分”自然而然地汇入了当时正引导人们对“连续统”及实变量函数进行持续研究（魏尔斯特拉斯、杜布瓦–雷蒙、汉克尔、迪尼）、并最终在康托尔那里开出集合论之花的思想潮流。黎曼赋予可积条件的形式，启发人们对函数在一个区间内的不连续点集考虑“测度”的观念；但几乎过了 30 年，才有人给出这一概念的一个富有成效且便于使用的定义。

这个方向上的最初尝试属于施托尔茨、哈纳克和康托尔（1884–1885 年）；前两人为了定义 R 的有界子集 E 的“测度”，考虑包含 E 的集合 \(F \supset E\)，即有限个区间之并，对每个 F 取相应区间长度之和，并把 E 的“测度”说成是这些数的下确界；而康托尔则从一开始就在空间 \(R^{n}\) 中工作，他对有界集 E 以及 \(\rho > 0\) 考虑 E 的邻域 \(V(\rho)\)（由到 E 的距离 \(\leq \rho\) 的点构成），并取 \(V(\rho)\) 的“体积”的下确界。⁶按照这个定义，一个集合的“测度”等于其闭包的“测度”，由此特别可知，两个没有公共点的集合之并的“测度”有可能严格小于这两个集合的“测度”之和。无疑是为了消除后一困难，佩亚诺{[}246 a{]}和若尔当（{[}174 c{]}，第 I 卷，第 28–31 页）在几年之后，对包含在一个方块 I 中的集合 A，在康托尔的“测度” \(\mu(A)\) 之外又引入它的“内测度” \(\mu(I) - \mu(I - A)\)，并把这两个数相等的那些集合 A（现称“可求积的”）称为“可测的”。这时两个没有公共点的可求积集 A、B 之并是可求积的，其“测度”等于 A 与 B 的“测度”之和；但有界开集不一定可求积，包含在一个有界区间中的有理数集亦然，这就使佩亚诺–若尔当的概念失去了大部分意义。

E. 博雷尔{[}32 a{]}的功绩在于能够看出以往各种定义的缺陷，并且看到了如何补救它们。自康托尔以来人们已经知道，R 中的每个开集 U 都是它的可数多个“分支”即两两不交的开区间之并；博雷尔不再试图用有限个区间把 U 围住而从“外部”逼近它，而是依靠上述结果，提议把 U 的各分支长度之和取作 U 的（当 U 有界时的）测度。随后他十分简略地\(^{7}\)描述了可以从开集出发、把可数并和“差” A - B 这两种运算无限迭代而得到的那个集合类（也称“博雷尔”集），并指出对这类集合可以定义一个具有完全可加性这一基本性质的测度：如果序列 \((A_{n})\) 由两两不交的博雷尔集组成，则它们之并（假定有界）的测度等于它们的测度之和。

这个定义将为分析开创一个新纪元：一方面，它与贝尔同时代的工作相配合，成为关于点集分类的一整系列拓扑性质研究的出发点（参见第 165 页）；而尤其重要的是，它将成为勒贝格在 20 世纪初头几年中所创造的积分概念推广的基础。

勒贝格在他的论文{[}196 a{]}中，首先精确化并发展了 E. 博雷尔的简略提示；他仿照佩亚诺–若尔当的方法，把有界集 \(A \subset \mathbf{R}\) 的“外测度”定义为包含 \(A\) 的开集的测度的下确界；然后，若 \(I\) 是包含 \(A\) 的一个区间，则 \(A\) 的“内测度”是 \(I\) 与 \(I - A\) 的外测度之差；这样就得到了“可测集”的概念，它与博雷尔最初的“构造性”定义的差别只在于添上了博雷尔意义下测度为零的集合的一个子集。这个定义立即推广到了空间 \(\mathbf{R}^n\)；而把定积分 \(\int_{a}^{b} f(t) dt\)（函数有界且 \(\geq 0\)）看作由曲线 \(y = f(x)\)、直线 \(x = a, x = b\) 与 \(y = 0\) 所围“面积”的那个旧观念，就这样向所有使上述集合之测度有定义的函数 \(f\) 提供了黎曼积分的一个直接推广。但是

勒贝格的独创性并不那么在于这一推广的思想\(^{8}\)，而在于他发现了关于在这样构想的积分中过渡到极限的基本定理，这个定理在他的工作中表现为测度完全可加性的一个推论；\(^{9}\)他立即认识到它的全部重要性，并把它作为基石，用于他早在 1904 年于著名的《积分论与原函数探求讲义》{[}196 c{]}中所给出的其理论的讲授性阐述。\(^{10}\)

我们在这里无法详细描述勒贝格的结果在无穷小演算古典问题的研究中所带来的不可胜数的进步。勒贝格本人在其论文中就已经把他的理论应用于把长度和面积这些古典概念推广到比通常曲线和曲面更一般的集合；关于这一理论半个世纪以来的可观发展，我们请读者参看 L. 切萨里最近的阐述{[}59{]}。还要提到应用于三角级数的那些成果，它们是勒贝格在其论文完成后几乎立即发展的{[}196 b{]}，并为这门理论开辟了新的视野，对其探索时至今日仍远未结束（参见{[}345{]}）。最后尤其是，空间 \(L^p\) 的定义和菲舍尔–里斯定理（{[}111{]}，{[}260 a and c{]}；参见第 213 页）揭示了积分这个新概念在泛函分析中可能起到的作用；这一作用只能随着这一概念后来的种种推广（我们马上就要谈到）而不断增长。

在此之前，我们要在勒贝格倾注心力最多的那些问题之一上稍作停留，即积分与原函数这两个概念之间的联系。随着黎曼所引入的积分的推广，自然产生的问题是：对连续函数有效的积分与原函数之间的古典对应，在更一般的情形下是否仍然成立。但很容易举出黎曼意义下可积的函数 f 的例子，使得 \(\int_{a}^{x}f(t)dt\) 在某些点没有导数（甚至也没有右导数或左导数）；

反过来，沃尔泰拉早在 1881 年就证明，一个函数 \(F(x)\) 可以在一个区间 \(I\) 中有有界的导数，却在 \(I\) 中（黎曼意义下）不可积。通过一种极为精细的分析（其中关于积分中过渡到极限的定理远远不够用），勒贝格成功地证明：若 \(f\) 在 \([a,b]\) 中（按他的意义）可积，则 \(F(x) = \int_{a}^{x} f(t) dt\) 几乎处处有等于 \(f(x)\) 的导数{[}196 c{]}。反之，若函数 \(g\) 在 \([a,b]\) 中可微且其导数 \(g' = f\) 有界，则 \(f\) 可积，并且我们有公式 \(g(x) - g(a) = \int_{a}^{x} f(t) dt\)。但勒贝格注意到，当 \(g'\) 不有界时问题要复杂得多；这时 \(g'\) 不一定可积，因此第一个问题便是刻画这样的连续函数 \(g\)：其 \(g'\) 几乎处处存在且可积。把自己限制在其诸导数\(^{11}\)之一处处有限的那些函数 \(g\) 上，勒贝格证明了这样的 \(g\) 必为有界变差函数。\(^{12}\)最后他建立了这最后一个结果的逆命题：有界变差函数 \(g\) 几乎处处有导数，且 \(g'\) 可积；但不再必然成立的是

\[
g (x) - g (a) = \int_ {a} ^ {x} g ^ {\prime} (t) d t;\tag{22.2}
\]

这个关系两边之差是一个有界变差函数，它不恒为常数且几乎处处导数为零（一个“奇异”函数）。剩下要刻画的是使关系式 (22.2) 成立的那些有界变差函数 g。勒贝格建立了：这些函数（维塔利称之为“绝对连续”函数，并对它们作了详细研究）正是具有下述性质的函数：g 在开集 U 中的全变差（即 g 在 U 的每个连通分支中的全变差之和）随 U 的测度趋于 0。

我们在下文将会看到，这些结果后来如何以一种弱化了的形式获得远为一般的适用范围。就其最初的形式而言，它们的应用领域仍相当有限，没有超出实变量函数“精细”理论的框架，对这一理论的后续发展我们在这里不作研究；我们只提一下当茹瓦及其追随者和后继者（佩龙、德拉瓦莱–普桑、辛钦、卢津、巴拿赫等）的深刻工作；读者可在 S. 萨克斯的书{[}268{]}中找到详细的阐述。

勒贝格理论所带来的实质性进步之一涉及多重积分。这个概念大约在 18 世纪中叶被引入，而且首先是作为“不定积分”的形式（通过与一元函数积分论的类比，\(\int\int f(x,y)dx dy\) 表示方程 \(\frac{\partial^{2}z}{\partial x\partial y}=f(x,y)\) 的一个解）；但早在 1770 年，欧拉就已经对展布在有界区域（由解析曲线弧围成）上的二重积分有了非常清楚的概念，并且正确地写出了用两个依次进行的简单积分来计算这种积分的公式（{[}108 a{]}，(1)，第 XVII 卷，第 289–315 页）。只要被积函数连续、积分区域也不太复杂，从“黎曼和”出发证明这个公式并不困难；但一旦想处理更一般的情形，黎曼的程序就遇到了严重的困难\((f(x,y)\) 可以在黎曼意义下可积，而当各简单积分按黎曼意义取时，\(\int dx\int f(x,y)dy\) 却没有意义）。当我们转到勒贝格的定义时，这些困难就消失了；勒贝格早在其论文中就证明了：当 \(f(x,y)\) 是有界的“贝尔函数”时，函数 \(y\mapsto f(x,y)\)（对一切 x）和 \(x\mapsto\int f(x,y)dy\) 也如此，并且我们有公式

\[
\int \int f (x, y) d x d y = \int d x \int f (x, y) d y\tag{22.3}
\]

（积分在矩形上取）。稍后，富比尼{[}121{]}为这一结果补充了一个重要的定理：若只假设 \(f\) 可积，则 \(x\) 的那些值（对它们而言 \(y \mapsto f(x, y)\) 不可积）构成一个零测集，这就使公式 (22.3) 得以推广到这种情形。

最后，在 1910 年{[}157 e{]}，勒贝格着手把他关于简单积分之导数的结果推广到多重积分。这样他就被引导着对一个函数 \(f\)（它在 \(\mathbf{R}^n\) 的所有紧子集上可积）联结上集合的函数 \(F(E) = \int_E f(\mathbf{x}) dx\)，它对 \(E\) 取遍 \(\mathbf{R}^n\) 的一切可积子集都有定义，推广了“不定积分”的概念；他借此机会注意到这个函数具有以下两条性质：\(1^\circ\) 它是完全可加的；\(2^\circ\) 它是“绝对连续的”，意思是 \(F(E)\) 随 \(E\) 的测度趋于 0。勒贝格这篇论文的实质部分在于证明这个命题的逆命题。\(^{13}\)但他没有止步于此，而是在同一方向上指出了推广有界变差函数概念的可能性，即考虑可测集的函数 \(F(E)\)，它完全可加，并且 \(\sum_n |F(E_n)|\) 对 \(E\) 分成可测子集 \(E_n\) 的每个可数分割都保持有界。而且，尽管他实际上只限于在 \(\mathbf{R}^n\) 的方块的集合上考虑这样的函数，但十分清楚，只需再走一步就可得到 J. 拉东将在 1913 年定义的一般的测度概念，它把勒贝格积分与斯蒂尔杰斯积分汇集于同一个综合之中，对后一种积分我们现在必须谈一谈。

1894 年，T. 斯蒂尔杰斯以《连分数研究》为题{[}297{]}发表了一篇极具独创性的论文，其中从表面看来十分狭窄的问题出发，以罕见的优雅陈述并解决了在解析函数论和实变量函数论中性质全新的问题。\(^{14}\)为了表示解析函数某个序列的极限，斯蒂尔杰斯（除其他事情外）被引导着在直线上引入正“质量分布”的概念，这个概念在物理科学中久已为人熟悉，但在此之前，它在数学中还从未在限制性假设（一般地，即在每个点存在连续变化的“密度”）之外被考虑过；他指出，给出这样一个分布等价于给出一个递增函数 \(\phi(x)\)，它对 x \textgreater{} 0 给出端点为 0 与 x 的区间中所含的总质量，对 x \textless{} 0 则给出改变符号后的这个质量，而 \(\phi\) 的不连续点对应于“集中在一点上”的质量。\(^{15}\)接着，斯蒂尔杰斯对区间 \([a, b]\) 中的这样一种质量分布构造“黎曼和” \(\sum_{i} f(\xi_{i})(\phi(x_{i+1}) - \phi(x_{i}))\)，并证明当 f 在 \([a, b]\) 中连续时，这些和趋于一个极限，他把它记作 \(\int_{a}^{b} f(x) d\phi(x)\)。由于只需要积分连续函数（甚至是可微函数），斯蒂尔杰斯没有再深入研究这个积分\(^{16}\)，在大约十年间这个概念似乎不曾引起注意。\(^{17}\)但 1909 年，F. 里斯{[}260 d{]}在解决阿达马几年前提出的一个问题（参见第 213 页）时，证明了斯蒂尔杰斯积分 \(f \mapsto \int_{a}^{b} fd\phi\) 是空间 \(\mathcal{C}(I)\)（在 \(I = [a, b] (\mathcal{C}(I))\) 上连续的数值函数的空间，赋有一致收敛拓扑\(^{18}\)）上最一般的连续线性泛函；这个结果的优雅与简洁几乎立即引起了它的几个推广。其中最成功的是 J. 拉东 1913 年的推广{[}256{]}：他把 F. 里斯的思想与勒贝格的思想结合起来，证明了如何用勒贝格的程序、从一个任意的“完全可加的集合函数”（定义在勒贝格测度意义下的可测集上）出发而不是从勒贝格测度出发，来定义一个积分。在这样定义的 \(R^{n}\) 上的“拉东测度”概念中，“有界变差”函数的概念被吸收了：把这样一个函数分解为两个递增函数之差，是测度分解为两个正测度之差的一个特例；同样，“以 \(\mu\) 为基的测度”对应于“绝对连续”函数的概念，而把任意测度分解为一个以 \(\mu\) 为基的测度和一个与 \(\mu\) 互外的测度，对应于有界变差函数分解为绝对连续函数与“奇异”函数之和的勒贝格分解。此外，拉东还证明了，关于 \(\mu\) 的“密度”——当一个测度以 \(\mu\) 为基时——在 \(\mu\) 是以勒贝格测度为基的测度的情形下仍然存在，他所用的是 F. 里斯先前的一个想法（后来被 J. 冯·诺伊曼等人重新提起并普及），即构造测度 \(\mu\) 在一个映射 \(\theta\)（由 \(R^{n}\) 到 R，其选取使 \(\theta(\mu)\) 是 R 上的勒贝格测度）下的象。

拉东的论文发表后几乎立即，弗雷歇就指出，这项工作中的几乎全部结果都可以推广到如下情形：“完全可加的集合函数”不定义在 \(R^{n}\) 的可测集上，而是定义在一个任意集合 E 的某些子集上（这些子集使得可数并与“差”这两种运算仍然给出该函数有定义的集合）。然而，把以 \(\mu\) 为基的测度表成 g. \(\mu\) 的形式这件事，在勒贝格和拉东那里依赖于本质上使用了 \(R^{n}\) 的拓扑的论证（而且我们已经看到，拉东的证明只有当 \(\mu\) 是以勒贝格测度为基的测度时才适用）；直到 1930 年，O. 尼科迪姆{[}235{]}才用直接的论证得到了一般形式的这个定理（几年后 J. 冯·诺伊曼借助空间 \(L^{2}\) 的性质（{[}324 g{]}，第 127–130 页）对它作了显著的简化）。

有了拉东的论文，积分的一般理论就可以认为在其主要轮廓上已经完成；作为其后的实质性收获，只能提到由丹尼尔{[}76 b{]}给出的测度无穷乘积的定义，以及博赫纳 1933 年给出的取值于巴拿赫空间的函数的积分{[}24 a{]}，后者是几年后由盖尔范德{[}125 a{]}、邓福德和佩蒂斯（{[}95{]}与{[}96{]}）所发展的更一般的“弱积分”研究的先声。但还剩下使新理论普及开来、使它成为普遍使用的数学工具的工作，因为在 1910 年前后，大多数数学家在“勒贝格积分”中仍然只看到一件精密度很高、操作细腻、只适用于极端精细和极端抽象的研究的工具。这就是卡拉泰奥多里在一本长期保持经典地位的著作{[}49{]}中所做的工作，这部书以大量独创性的评注充实了拉东的理论。

但也正是在这部书中，一直处于勒贝格关注之首的积分概念（他的论文{[}196 a{]}和他关于这些问题的主要著作{[}196 c{]}的标题就足以表明这一点），第一次把地盘让给了测度概念，而后者对勒贝格来说（一如在他之前对若尔当那样）只是一种辅助的技术手段。观点上的这一变化，在卡拉泰奥多里那里无疑是由于他似乎赋予了“p 维测度”过大的重要性。\(^{19}\)从那时起，论述积分的作者们就在这两种观点之间分成了两派，而且不免卷入了一些引起大量墨水（如果不是大量鲜血的话）之争的论战。一派追随着卡拉泰奥多里；在他们越来越抽象、越来越公理化的阐述中，测度连同它所允许的一切精细加工，不仅扮演着主导角色，而且趋于失去与拓扑结构的联系，而在测度所涉及的大多数问题中，它实际上是与这些结构相连的。另一些阐述则或多或少紧随 W. H. 杨早在 1911 年就已指出、其后由丹尼尔发展的方法，那篇指出该方法的论文可惜很少受人注意{[}339 b{]}。第一种阐述在处理勒贝格积分时，从假定已知的、具有紧支集的连续函数的“柯西积分”出发，依次定义下半连续函数的上积分，然后定义任意数值函数的上积分，由此用纯“泛函”的手段得到可积函数的一个定义，它是仿照勒贝格对集合的定义作出的。丹尼尔在 1918 年（{[}76 a{]}；参见{[}207{]}）把这一阐述（有一些变化）推广到定义在任意集合上的函数；在同一思想脉络中（并与盖尔范德之前谱论中所用的方法密切相关），还必须指出 F. 里斯的论文{[}260 h{]}，它以优雅而简洁的形式给出了序空间理论中在积分理论里起作用的那些结果。

自 1920 年以来源于积分理论的进步，与其说须在读起来多少令人愉快的阐述性著作中去寻找——它们的实质内容不可能有多大的变化——不如说须在应用方面去寻找：概率论（昔日曾是猜测与悖论的托词，自其

被科尔莫哥洛夫公理化{[}183{]}之后，已成为积分论的一个分支，而且是一个有着自己的方法和问题的自主分支）；遍历理论；谱理论和调和分析，因为哈尔发现以他的名字命名的测度以及这一发现所引发的思想运动，已使积分成为群论中最重要的工具之一。
% semantic-fragments: legacy-input-seam
\relax{}
